%=====================================================================
\section{The order $t^0$ coefficient}\label{sec:order0}
%=====================================================================
Notation is that of Section~\ref{sec:radial}. Put $B_1(u)=u-\frac12$ on $[0,1)$ and $B_2(u)=\{u\}^2-\{u\}+\frac16$ (periodic, even, continuous).

\begin{lemma}\label{lem:EMline}
Let $\phi:\R\to\C$ be smooth on $(-\infty,0]$ and on $[0,\infty)$ separately, with one-sided derivatives of order $\le3$ bounded by $c_1e^{-c_2u^2}$. For $N\ge1$ and
$\rho\in\R$ let $\Sigma(h)=\sum_{u\in\rho+N\Z}{}'\phi(hu)$, the prime giving weight $\frac12(\phi(0^+)+\phi(0^-))$ to $u=0$. Then as $h\to0^+$,
\[ \Sigma(h)=\frac1{Nh}\int_\R\phi-\bar B_1(\rho/N)\big(\phi(0^+)-\phi(0^-)\big)-\frac{Nh}2B_2(\rho/N)\big(\phi'(0^+)-\phi'(0^-)\big)+O(h^2). \]
\end{lemma}
\begin{proof}
Apply Lemma~\ref{lem:EM1} (third line, remainder $O(N^2\int|h'''|)$ which is $O(h^2)$ after the scaling $u\mapsto hu$) to $u\mapsto\phi(hu)$ on $u\ge0$ and to
$u\mapsto\phi(-hu)$ on $u>0$ with offset $\{-\rho/N\}$; use $B_1(1-y)=-B_1(y)$ and $B_2(1-y)=B_2(y)$ and add the half-weight at $u=0$.
\end{proof}

\begin{theorem}\label{thm:A2odd}
For every companion and rational $x$, $A_2(S^o)=0$.
\end{theorem}
\begin{proof}
By Theorem~\ref{thm:radtriv}, $A_2(S^o)=-\sum_c\kappa_cA_2(R_c)$. In the proof of Corollary~\ref{cor:A1odd}, the sum along the wall direction is
$\varphi(\rho)\sqrt{\pi/(tQ(d))}e^{ts_c^2/4}$ up to $O(t^K)$, and the transverse sum is $\sum_{a\in\rho+N\Z}h(\sqrt ta)$ with $h$ odd. Then $h(0^+)-h(0^-)=-2$ and,
$h'$ being even, $h'(0^+)-h'(0^-)=0$. Lemma~\ref{lem:EMline} gives $2\bar B_1(\rho/N)+O(t)$, so $R_c$ has no $t^0$ term.
\end{proof}

\begin{theorem}\label{thm:A2even}
For every companion and rational $x$, with the coset classes $r$ of $\alpha+\Z^2$ modulo a common period $N$ of the phase, $f_r$ the phase on $r$ and $(\rho_1,\rho_2)$ the
coordinates of a representative,
\[ \tfrac12A_2(S^E)=\sum_r f_r\Big[\bar B_1(\rho_1/N)\bar B_1(\rho_2/N)+\tfrac12B_2(\rho_1/N)+\tfrac74B_2(\rho_2/N)\Big]=:A_2^{\rm corner}. \]
\end{theorem}
\begin{proof}
Write $E=\chi_{++}+\chi_{--}$ as in Theorem~\ref{thm:evenA1} and $h(w)=e^{-Q(w)}$, a Gaussian on the closed first quadrant. On a class $r$ the $\chi_{++}$-sum is
$\sum'_{\nu_1}\sum'_{\nu_2}h(\sqrt t\nu)$ over $\nu_i\in\rho_i+N\Z$, $\nu_i\ge0$. Apply Lemma~\ref{lem:EMline} with $\phi=h(\sqrt t\nu_1,\cdot)\mathbf 1_{\ge0}$ in $\nu_2$, then in
$\nu_1$ to each resulting term; the remainders are uniform because $h$ and its derivatives are Gaussian on the quadrant. Collecting the $t^0$ terms gives
$\bar B_1(\rho_1/N)\bar B_1(\rho_2/N)h(0)-\frac12B_2(\rho_1/N)\int_0^\infty\partial_1h(0,w)\,dw-\frac12B_2(\rho_2/N)\int_0^\infty\partial_2h(v,0)\,dv$. Here
$\partial_1h(0,w)=-7we^{-\frac72w^2}$ and $\partial_2h(v,0)=-7ve^{-v^2}$, with integrals $-1$ and $-\frac72$. For $\chi_{--}$ substitute $\nu=-\mu$: the same expression at
$(-\rho_1,-\rho_2)$, which is equal since $\bar B_1\bar B_1$ and $B_2$ are even. Hence $A_2(S^E)$ on the class is twice the bracket.
\end{proof}

For the line terms let $A_2^{\rm lines}$ be the sum over the lines of Proposition~\ref{prop:lines} of
$w_+\big[-\mean(f)\frac{l}{2Q(u)}-\sum_{\rho<N}f(\rho)B_1(\rho/N)\big]+w_-\big[\mean(f)\frac{l}{2Q(u)}-\sum_{\rho<N}f(-\rho)B_1(\rho/N)-f(0)\big]+\sum_b\delta(b)f(b)$
(with $l=B(\nu_0,u)$, $f(b)$ the phase at $\nu_0+bu$, $\delta$ the finite part of the weight), plus the point terms $h_1h_2$; this is the $t^0$ term of
Theorem~\ref{thm:expansion}'s ray formula applied to each half-line.
\begin{theorem}\label{thm:A2exact}
$A_2=A_2^{\rm lines}+A_2^{\rm corner}$. In particular $A_2-A_2^{\rm lines}$ is $P$-even, and the completion terms contribute nothing at order $t^0$.
\end{theorem}
\begin{proof}
$S^W=S^{W-p}+S^o+\frac12S^E$; apply Theorems~\ref{thm:A2odd} and~\ref{thm:A2even} and the ray formula.
\end{proof}

\begin{proposition}\label{prop:noaffine}
Let $a_0,a_2:\Q^\times\to\C^{105}$ with $a_0(-1/x)=x^{-1}35^{-1/2}Ma_0(x)$. If $a_2(-1/x)=x\,35^{-1/2}Ma_2(x)+c\,Ma_0(x)$ for all $x\in\Q^\times$ and a constant $c$, then
$Pa_2=-a_2$. Consequently $A_2-A_2^{\rm lines}$ satisfies such a law on $\Q^\times$ only if it vanishes.
\end{proposition}
\begin{proof}
Apply the law at $-1/x$ and use $M^2=35P$: $a_2(x)=-Pa_2(x)-\frac cx35^{-1/2}M^2a_0(x)+cM\,x^{-1}35^{-1/2}Ma_0(x)=-Pa_2(x)$. The second statement follows from
Theorem~\ref{thm:A2exact}.
\end{proof}
\begin{remark}
If a vector $F$ satisfied $F(-1/\tau)=\tau35^{-1/2}MF(\tau)$, expanding at $\tau=x+i\varepsilon$ with $w=\varepsilon/(x^2+ix\varepsilon)$ gives the order-$t^0$ relation
$a_2(-1/x)=x35^{-1/2}Ma_2(x)$: the terms in $a_0$ from $\tau=x+i\varepsilon$ and from $1/w=x^2/\varepsilon+ix$ cancel, so the constant of Proposition~\ref{prop:noaffine}
produced by a weight-one law is $c=0$.
\end{remark}
\begin{catbremark}
On the $216$ points $x\in(0,3]$ the law of the remark fails for $A_2-A_2^{\rm lines}$ with absolute residual $54.4$ (relative $0.059$); the three parts of $A_2^{\rm corner}$
satisfy no law $X(-1/x)=x^w35^{-1/2}MX(x)$, $w\in\{-1,0,1\}$ (Appendix~\ref{app:catB}, item~\ref{catB:o19}).
\end{catbremark}
